% Propietario conservado: /Users/ruben/Documents/New project/output/EXTREMA_MEDIA_RAZON_RECIPROCIDAD_ES_EN_20260918/ARTICULO/ES/source/nuclear/gaussianas.tex
\chapter{Covarianza, forma normal y espectro gaussiano}
\label{x_reciprocidad:app:gaussianas-espectro}

La sección positiva de acción \(\hbar_a\) y la realización canónica
anterior determinan la normalización empleada aquí. Se demuestra la
correspondencia entre covarianza, espectro, pureza y entropía para un
número finito de modos; el espacio de Fock conserva todas sus ocupaciones.
La reducción simpléctica y las fórmulas gaussianas son resultados
establecidos que se reproducen para hacer autónoma su aplicación al
transporte de memoria.

\section{Coordenadas canónicas y función característica}
En \(L^2(\RR^d)\), sean \(\widehat q_j=\widehat Q_j/\sqrt{\hbar_a}\),
\(\widehat p_j=\widehat P_j/\sqrt{\hbar_a}\) y
\(\widehat z=(\widehat q_1,\widehat p_1,\ldots,\widehat q_d,\widehat p_d)\).
Sobre Schwartz,
\[
 [\widehat z_j,\widehat z_k]=iJ_{jk}I,\qquad
 J=\bigoplus_{j=1}^d\begin{pmatrix}0&1\\-1&0\end{pmatrix}.
\]
Esta convención fija el signo de la forma canónica. Los operadores de
Weyl \(\mathcal W(\xi)=\exp(i\xi^{\mathsf T}\widehat z)\) se definen
por las clausuras autoadjuntas de las combinaciones lineales indicadas.
Para un estado centrado \(\rho\) de segundos momentos finitos, la
covarianza normalizada y su función característica son
\[
 W_{jk}=\Tr\rho(\widehat z_j\widehat z_k+\widehat z_k\widehat z_j),
 \qquad \chi_\rho(\xi)=\Tr(\rho\mathcal W(\xi)),\qquad
 V=\frac{\hbar_a}{2}W.
\]
El estado es gaussiano centrado cuando
\begin{equation}
 \chi_\rho(\xi)=\exp(-\xi^{\mathsf T}W\xi/4).
 \label{x_reciprocidad:app:eq:caracteristica-gaussiana}
\end{equation}
La esperanza de \((c^{\mathsf T}\widehat z)^\dagger
(c^{\mathsf T}\widehat z)\) es \(c^\dagger(W+iJ)c/2\); así
\(W+iJ\geq0\). Además \(W>0\): si \(v\) es real y
\(v^{\mathsf T}Wv=0\), la positividad implica \((W+iJ)v=0\),
por tanto \(Jv=0\) y \(v=0\).

\begin{lemma}[Unicidad de la función característica]
\label{x_reciprocidad:app:lem:unicidad-weyl}
Dos operadores de clase traza con la misma función característica son iguales.
\end{lemma}
\begin{proof}
Reagrupando posiciones y momentos como \(\xi=(u,v)\),
\[
 (\mathcal W(u,v)f)(x)=e^{iu\cdot(x+v/2)}f(x+v).
\]
Para un operador de rango finito con vectores de Schwartz y núcleo \(K_A\),
\[
 \Tr(A\mathcal W(u,v))=\int K_A(x+v/2,x-v/2)e^{iu\cdot x}\,dx.
\]
Plancherel y el cambio de variables de jacobiano absoluto uno dan
\[
 \|A\|_{\rm HS}^2=(2\pi)^{-d}\int_{\RR^{2d}}
 |\Tr(A\mathcal W(\xi))|^2\,d\xi.
\]
Los operadores anteriores son densos en clase traza. La convergencia
en esa norma implica convergencia en Hilbert--Schmidt y convergencia
uniforme de las funciones características, porque su diferencia está
acotada por \(\|A-B\|_1\). La igualdad se extiende en \(L^2\).
Aplicarla a la diferencia de los operadores prueba la afirmación.
\end{proof}

\section{Reducción simpléctica de la forma positiva}
\begin{theorem}[Forma normal de Williamson]
\label{x_reciprocidad:app:thm:williamson}
Para una matriz real simétrica \(W>0\) de orden \(2d\), existen
\(S\in\operatorname{Sp}(2d,\RR)\) y \(\nu_j>0\) tales que
\[
 W=SDS^{\mathsf T},\qquad D=\bigoplus_j\nu_jI_2.
\]
Los autovalores de \(JW\) son \(\pm i\nu_j\). La condición
\(W+iJ\geq0\) equivale a \(\nu_j\geq1\) para todo \(j\).
\end{theorem}
\begin{proof}
La matriz \(B=W^{1/2}JW^{1/2}\) es antisimétrica e invertible.
Si \(-B^2u=\nu^2u\) y \(\|u\|=1\), el vector
\(v=-Bu/\nu\) es unitario y ortogonal a \(u\); cumple
\(Bu=-\nu v\), \(Bv=\nu u\). El complemento de esta pareja es
invariante. La iteración produce una matriz ortogonal \(O\) con
\(O^{\mathsf T}BO=JD\). Definiendo \(S=W^{1/2}OD^{-1/2}\),
\[
 S^{\mathsf T}JS=J,\qquad SDS^{\mathsf T}=W.
\]
La semejanza \(W^{1/2}(JW)W^{-1/2}=B\) determina los valores
\(\nu_j\), con multiplicidad. La congruencia por \(S^{-1}\)
transforma \(W+iJ\) en \(D+iJ\); cada bloque tiene autovalores
\(\nu_j+1\) y \(\nu_j-1\).
\end{proof}

\begin{lemma}[Implementación unitaria en un número finito de modos]
\label{x_reciprocidad:app:lem:implementacion-simplectica}
Toda matriz simpléctica \(S\) admite un unitario \(U_S\) que satisface
\(U_S^\dagger\mathcal W(\xi)U_S=\mathcal W(S^{\mathsf T}\xi)\).
La conjugación \(\rho\mapsto U_S\rho U_S^\dagger\) transporta
la covarianza por \(W\mapsto SWS^{\mathsf T}\).
\end{lemma}
\begin{proof}
En la descomposición polar \(S=OP\), la identidad
\(J^{-1}(S^{\mathsf T}S)J=(S^{\mathsf T}S)^{-1}\), aplicada por
cálculo espectral a la raíz positiva, prueba \(PJP=J\).
Entonces \(O\) es ortogonal simpléctica. Los autoespacios de \(P\)
se emparejan mediante \(J\) con autovalores \(\lambda,\lambda^{-1}\);
el de autovalor uno es \(J\)-invariante. Por ello una matriz
ortogonal simpléctica \(K\) reduce \(P\) a
\(\bigoplus_j\operatorname{diag}(e^{r_j},e^{-r_j})\).

La dilatación se implementa por
\[
 (\mathcal D_rf)(x)=e^{-\frac12\sum_jr_j}
 f(e^{-r_1}x_1,\ldots,e^{-r_d}x_d).
\]
Un cambio de variables prueba su unitariedad y la conjugación de
\((\widehat q_j,\widehat p_j)\) a
\((e^{r_j}\widehat q_j,e^{-r_j}\widehat p_j)\).
Una transformación ortogonal simpléctica conmuta con \(J\) y
corresponde a una matriz \(u\in U(d)\) sobre
\(a_j=(\widehat q_j+i\widehat p_j)/\sqrt2\). Su segunda
cuantización \(\Gamma(u)=\bigoplus_{n\geq0}u^{\otimes_s n}\)
es unitaria en \(\mathcal F_s(\CC^d)\), y la identidad sobre
productos simétricos finitos da la transformación de los \(a_j\).

La base de Hermite identifica este espacio con \(L^2(\RR^d)\).
Para justificar su completitud, si \(f\) es ortogonal a todos los
polinomios multiplicados por \(e^{-|x|^2/2}\), la función entera
\(F(z)=\int f(x)e^{-|x|^2/2}e^{z\cdot x}\,dx\) tiene nulas
todas sus derivadas en cero. Así \(F=0\); restringiendo a
\(z=i\xi\), la unicidad de Fourier da \(f=0\). Las
normalizaciones y la ortogonalidad resultan de creación y aniquilación.
La composición de las implementaciones anteriores produce \(U_S\).
La igualdad de Weyl implica
\(\chi_{U_S\rho U_S^\dagger}(\xi)=\chi_\rho(S^{\mathsf T}\xi)\),
que demuestra la ley de covarianza.
\end{proof}

\section{Cálculo del espectro de la densidad isotrópica}
En un modo, los estados coherentes normalizados son
\[
 |z\rangle=e^{-|z|^2/2}\sum_{n\geq0}\frac{z^n}{\sqrt{n!}}|n\rangle.
\]
Su función de posición,
\(\pi^{-1/4}\exp[-(x-q_0)^2/2+ip_0(x-q_0/2)]\), donde
\(q_0=\sqrt2\Re z\), \(p_0=\sqrt2\Im z\), da por integración
\[
 \langle z|\mathcal W(u,v)|z\rangle
 =e^{-(u^2+v^2)/4}e^{i(uq_0+vp_0)}.
\]
Para \(s>0\), la integral
\[
 \rho^{(s)}=\frac1{\pi s}\int_{\CC} e^{-|z|^2/s}
 |z\rangle\langle z|\,d^2z
\]
converge en clase traza, es positiva y tiene traza uno. La integración
angular anula sus elementos fuera de la diagonal de ocupación; la
integración radial da
\[
 \langle n|\rho^{(s)}|n\rangle=\frac{s^n}{(1+s)^{n+1}},\qquad
 \chi_{\rho^{(s)}}(u,v)=e^{-(2s+1)(u^2+v^2)/4}.
\]
Con \(\nu=2s+1\), el único estado gaussiano de covarianza
\(\nu I_2\) es, por el lema~\ref{x_reciprocidad:app:lem:unicidad-weyl},
\begin{equation}
 \rho_\nu=(1-t)\sum_{n\geq0}t^n|n\rangle\langle n|,
 \qquad t=\frac{\nu-1}{\nu+1}.
 \label{x_reciprocidad:app:eq:estado-geometrico}
\end{equation}
El caso \(\nu=1\) es el proyector sobre el vacío.

\begin{theorem}[Espectro, pureza y entropía]
\label{x_reciprocidad:app:thm:espectro-gaussiano}
Un estado gaussiano centrado con covarianza \(W\) y valores
simplécticos \(\nu_1,\ldots,\nu_d\) es unitariamente equivalente
a \(\bigotimes_j\rho_{\nu_j}\). Sus autovalores son
\[
 \lambda_{\symbf{n}}=\prod_j(1-t_j)t_j^{n_j},\qquad
 t_j=\frac{\nu_j-1}{\nu_j+1},\quad \symbf{n}\in\mathbb N^d.
\]
Con \(0\log0=0\),
\begin{align}
 \Tr\rho^2&=\prod_j\nu_j^{-1}=(\det W)^{-1/2},
 \label{x_reciprocidad:app:eq:pureza-general}\\
 -\Tr(\rho\log\rho)&=\sum_j
 \left[\frac{\nu_j+1}{2}\log\frac{\nu_j+1}{2}
 -\frac{\nu_j-1}{2}\log\frac{\nu_j-1}{2}\right].
 \label{x_reciprocidad:app:eq:entropia-general}
\end{align}
\end{theorem}
\begin{proof}
Williamson da \(W=SDS^{\mathsf T}\). El producto de las densidades
isotrópicas tiene función característica \(e^{-\xi^{\mathsf T}D\xi/4}\).
La implementación unitaria y la unicidad de Weyl prueban la
equivalencia y la fórmula espectral. En un modo,
\[
 \sum_{n\geq0}(1-t)^2t^{2n}=\frac{1-t}{1+t}=\nu^{-1},
\]
y
\[
 -\sum_{n\geq0}(1-t)t^n\log((1-t)t^n)
 =-\log(1-t)-\frac{t}{1-t}\log t.
\]
La sustitución \(t=(\nu-1)/(\nu+1)\) prueba las fórmulas; el
límite en \(t=0\) es cero. En el producto finito, Tonelli justifica
la suma de los términos de entropía no negativos, y
\(\det W=\prod_j\nu_j^2\) da la expresión determinantal.
\end{proof}

Si el estado es la reducción de un vector puro bipartito, sus
autovalores son los cuadrados de los coeficientes de Schmidt.
La descomposición singular del operador de Hilbert--Schmidt definido
por ese vector demuestra la igualdad de espectros no nulos de ambas
reducciones. La entropía calculada es entonces la de entrelazamiento.

\section{Preparación y transporte de las dos covarianzas}
La aplicación posterior utiliza dos entradas gaussianas puras idénticas.
En general, para \(M\in\operatorname{Sp}(2d,\RR)\), su covarianza es
\(V_0=(\hbar_a/2)MM^{\mathsf T}\), y la entrada conjunta tiene
covarianza \(V_0\oplus V_0\). El transporte físico se obtiene
conjugando \(U_H\) por \(M\oplus M\). En la carta equilibrada,
la covarianza normalizada de entrada es \(I_{4d}\); con
\[
 T=\frac{8I+H}{9},\qquad D=\frac{\sqrt8(H-I)}9,
\]
la primera reducción tiene
\[
 W=TT^{\mathsf T}+DD^{\mathsf T}=\frac{8I+HH^{\mathsf T}}9.
\]
Su función característica se obtiene anulando los argumentos del
segundo grupo de modos. Por ello permanece gaussiana y los teoremas
anteriores calculan todo su espectro, sin truncar el registro de ocupación.
